\section{A common four-site chain and six-coordinate boundary}
\label{sec:chain}

The remaining proofs use squared adjacency matrices. For a cap $c$,
positive definiteness of $cI-A^2$ is equivalent to $\rho(A)^2<c$.
Throughout, $\norm{\cdot}_2$ is the Euclidean operator norm and
$\norm{\cdot}_F$ is the Frobenius norm.
The required local output is a positive six-coordinate boundary matrix
with an explicit error bound after any sufficiently long interior chain.
We derive one block system and certify it separately at $c=\czero$ and
$c=\cone$; the second energy uses its own finite premises.

For the one-cell word $w_j$, put $h=8j+2=4\ell+2$, where $\ell=2j$.
Initially take positive holonomy and $h\ge18$. Partition the vertices into
\[
 V_0=\{0,1\},\qquad
 V_s=\{2+4(s-1),\ldots,5+4(s-1)\}\quad(1\le s\le\ell).
\]
The normalized matrix $M_h(c)=cI-A_h(w_j,+1)^2$ has bulk diagonal and
successive coupling blocks
\begin{gather}
 D_c=\begin{pmatrix}
 c-4&0&-1&0\\0&c-4&0&-1\\-1&0&c-4&0\\0&-1&0&c-4
 \end{pmatrix},\label{eq:D}\\
 E_+=\begin{pmatrix}-1&0&0&0\\0&1&0&0\\-1&2&1&0\\2&-1&0&-1\end{pmatrix},
 \qquad
 E_-=\begin{pmatrix}-1&0&0&0\\0&1&0&0\\-1&-2&1&0\\-2&-1&0&-1\end{pmatrix}.
 \label{eq:E}
\end{gather}
Here $M_h[V_s,V_{s+1}]=E_+$ for odd $s$ and $E_-$ for even $s$.
The boundary blocks are
\begin{gather}
 G_0=(c-4)I_2,\qquad H_0=D_c,\qquad
 R_0=\begin{pmatrix}-1&-2&1&0\\-2&-1&0&-1\end{pmatrix},\notag\\
 C_0=\begin{pmatrix}-1&0&-1&0\\0&-1&0&-1\end{pmatrix},\qquad
 W_0=\begin{pmatrix}0&0&-1&0\\0&0&0&1\\0&0&0&0\\0&0&0&0\end{pmatrix}.
 \label{eq:boundary-data}
\end{gather}
Specifically, $R_0=M_h[V_0,V_1]$, $C_0=M_h[V_0,V_\ell]$, and
$W_0=M_h[V_1,V_\ell]$. All other nonadjacent bulk blocks vanish.
These identities follow from the complete two-walk expansion
\begin{align}
 (A^2)_{i,i}&=4,& (A^2)_{i,i+1}&=\tau_{i-1}+\tau_i,&
 (A^2)_{i,i+2}&=1,\notag\\
 (A^2)_{i,i+3}&=\tau_i+\tau_{i+1},&
 (A^2)_{i,i+4}&=\tau_i\tau_{i+2},&&
 \label{eq:two-walks}
\end{align}
their transposes, and zero entries at all other cyclic offsets. At
$h\ge18$ these offsets and the stated block locations do not collide.
The exceptional order $h=10$ will be checked directly.

Use zero-based elimination indices: let $X_0=D_c$, with $E_a=E_+$ for
even $a$ and $E_a=E_-$ for odd $a$. Eliminating the current four-site
pivot gives the length-independent open recurrence
\begin{align}
 X_{a+1}&=D_c-E_a^TX_a^{-1}E_a,&
 R_{a+1}&=-R_aX_a^{-1}E_a,\notag\\
 W_{a+1}&=-E_a^TX_a^{-1}W_a,&
 G_{a+1}&=G_a-R_aX_a^{-1}R_a^T,\label{eq:recurrence}\\
 H_{a+1}&=H_a-W_a^TX_a^{-1}W_a,&
 C_{a+1}&=C_a-R_aX_a^{-1}W_a.\notag
\end{align}
The final pivot has index $p=\ell-2$, which is even. At that step the
right coupling is $W_p+E_+$, so the retained six-coordinate core is
\begin{equation}
 S_\ell(c)=\begin{pmatrix}
 G_p-R_pX_p^{-1}R_p^T&C_p-R_pX_p^{-1}(W_p+E_+)\\
 *&H_p-(W_p+E_+)^TX_p^{-1}(W_p+E_+)
 \end{pmatrix}.
 \label{eq:single-core}
\end{equation}
The lower-left block is the transpose of the upper-right block.
Repeated Schur congruence proves $M_h(c)\succ0$ once every eliminated
pivot and this core are positive. The subscript $\ell$ counts four-site
blocks, not vertices.

\subsection{The two finite rational certificates}

Write $F_\pm(X)=D_c-E_\pm^TX^{-1}E_\pm$ and
$\Phi=F_-\circ F_+$. At the even entrance $J$ set
\[
 Z=X_J=\Phi^{J/2}(D_c),\qquad Y=F_+(Z),\qquad
 T_Z=Z^{-1}E_+Y^{-1}E_-.
\]
Both energies use the fixed rational weights
\begin{align}
 P&=\frac1{10000}\begin{pmatrix}
10766&87&19&974\\87&12664&148&-2418\\19&148&10093&-25\\974&-2418&-25&14009
\end{pmatrix},\notag\\
 Q&=\frac1{10000}\begin{pmatrix}
11503&614&990&-1101\\614&10470&15&113\\990&15&12299&-2632\\-1101&113&-2632&13260
\end{pmatrix}.
\label{eq:weights}
\end{align}
The numerical constants are exact rational numbers:
\begin{center}
\begin{tabular}{lcc}
\toprule
Quantity & $c=\czero$ & $c=\cone$\\
\midrule
Entrance $J$ & $24$ & $48$\\
Riccati radius $\delta$ & $10^{-10}$ & $10^{-18}$\\
Response square bound $\zeta$ & $10^{-10}$ & $10^{-20}$\\
Center transfer factor $\kappa$ & $2/5$ & $1/2$\\
Response contraction $q$ & $2/3$ & $3/4$\\
Riccati contraction $\theta=q^2$ & $4/9$ & $9/16$\\
Even response bound $a$ & $1/30000$ & $1/9000000000$\\
Pair response bound $b=12a$ & $1/2500$ & $1/750000000$\\
Seed margin $\gamma$ & $1/50$ & $10^{-6}$\\
Tail tolerance $\epsilon$ & $1/500$ & $10^{-8}$\\
Seed graph order $4(J+2)+2$ & $106$ & $202$\\
\bottomrule
\end{tabular}
\end{center}

\begin{lemma}[Exact finite premises]\label{lem:finite-premises}
For each column of the table, the rational recurrence satisfies:
\begin{enumerate}
\item $X_0,\ldots,X_J\succ0$ and $Z,Y\succ I_4/2$;
\item $(9/10)I_4\prec P,Q\prec2I_4$;
\item $T_Z^TPT_Z\prec\kappa P$ and $T_ZQT_Z^T\prec\kappa Q$;
\item $\norm{\Phi(Z)-Z}_F<\delta/40$;
\item $R_JQR_J^T\prec\zeta I_2$ and $W_J^TQW_J\prec\zeta I_4$;
\item $S_{J+2}(c)\succ\gamma I_6$.
\end{enumerate}
\end{lemma}
\begin{proof}
All the matrices are specified by the rational initial data and at most
$J+2$ recurrence steps. The accompanying exact verifiers evaluate them over
$\mathbb Q$. For each positive-definiteness assertion they compute scalar
Schur pivots: if a symmetric matrix is partitioned as
$\left(\begin{smallmatrix}d&v^T\\v&K\end{smallmatrix}\right)$, they require
$d>0$ and recurse on $K-vv^T/d$. Every pivot is a strictly positive
rational number. The residual check is the exact comparison
\[
 \sum_{u,v}(\Phi(Z)-Z)_{uv}^2<(\delta/40)^2.
\]
The verifiers reconstruct the responses, and the certificates store the
exact centers, seeds and positive pivots; the $\cone$ certificate also
stores both entrance responses. An independent implementation constructs the full signed graph and
uses scalar elimination, without generating the entrance through the
four-site recurrence. It recovers the entrance after $4J$ scalar pivots,
subtracts the physical terminal $E_+$ coupling to isolate $W_J$, and
checks the same inequalities. Both energies pass these rational checks.
The implementation and certificate locations are given in
Section~\ref{sec:verification}.
\end{proof}

This lemma is the finite computational part of the uniform argument. No
floating eigenvalue or limiting-matrix approximation is used in its
acceptance tests. The seed margins refer to the normalized six-coordinate
cores in \eqref{eq:single-core}.
